%%
%% CONCEPTS
%%
\begin{slide}
  \begin{itemize}
    \item{Objects\,:}\\
      espace, temps, lois.
    \item{Transformations\,:}\\
      accumulations, diff\'erences.
    \item{Visualisations\,:}\\
      photographiques ou cin\'ematographiques, \\
      compl\`etes ou r\'eduite, \\
      solide ou transparente, \\
      projet\'ees, tranch\'ees, \'eclat\'ees ou coutour\'ees.
  \end{itemize}
\end{slide}

%%
%% ANNEAL 2K BITMAP
%%
\begin{slide}
  \mySfigure{full}{dump}%
  {anneal-g128-s11}{\arule{anneal} \asquare{2048} \agen{128}}
\end{slide}

\begin{mynote}
  ANNEAL 2K BITMAP

  Une photo directe, sans r\'eduction d'information.
  
  La r\'esolution du p\'eriph\'erique de sortie fixe une limite \`a la
  visualisation pratique (sans perte d'information) des configuration
  cellulaire.
\end{mynote}

%%
%% ANNEAL SEQUENCE
%%
\myslide{S\'equence de bitmaps}%
{%
  \mySfigure{full}{dump}%
  {anneal-03}{\arule{Anneal} \asquare{256} \agencnt{559}}}
{%
  ANNEAL SEQUENCE

  L'application it\'erative de la r\`egle sur la configuration initiale
  apporte la dimension temporelle au espace du calcul cellulaires.

  Et un probl\`eme de visualisation du \`a la taille en bits des
  trajectoires ou espace-temps.

  Sur cette s\'equence, l'intervalle de temps s\'eparant deux espace est
  incr\'ement\'e apr\`es chaque \'echantillonage.

  Une des solution \`a ce probl\`eme est la r\'ealisation de s\'equence anim\'e.

  J'ais r\'ealis\'e des visualisation cin\'ematographiques, mais je me
  limiterais dans le cadre de cette expos\'e au visualisations
  photographique.

  Je signale simplement les d\'ebits obtenu Sur un 486DX2 (\`a peu pr\`es 2
  fois une SparcStation ELC) pour \arule{Life} et \arule{Anneal}.

  Pour le calcul\,: 15 Hz.\\
  Pour l'animation (copy de fichier \`a fen\`etre)\,: 30 Hz.\\
  Pour la visualisation en temps r\'eel\,: 10 Hz.\\
  Pour \arule{Qmean} on obtient calcul\,:\\
  12 Hz, animation\,: 8 Hz, combin\'e\,: 5 Hz.}

%%
%% LIFE SUPERLUMINIC GLIDER
%%
\myslide{Glider Superluminique}
{\mySfigure{full}{dump}{superglider}{\arule{Life} \acube{128}}}
{Cette s\'erie d'espaces de configurations cellulaires est une s\'equence
  de tranches d'espace-temps non align\'ees sur l'axe du temps.}

%%
%% QMEAN BITMAPS SERIE
%%
\begin{slide}
  \mySfigure{full}{dump}{qmean-9-128}%
  {\arule{Qmean} \asquare{512} \agen{128} \aplane{3}}
\end{slide}

\begin{slide}
  \mySfigure{full}{dump}{qmean-9-512}%
  {\arule{Qmean} \asquare{512} \agen{512} \aplane{3}}
\end{slide}

\begin{slide}
  \mySfigure{full}{dump}{anneal-02}%
  {Deux configurations de la r\`egle~\arule{anneal}}%
\end{slide}

%%
%%
%%

\begin{slide}
  \mySfigure{full}{dump}{anneal-01}%
  {Les deux configurations de la figure pr\'ecedentes non retouch\'ee}%
\end{slide}

%%
%%
%%

%%
%% ANNEAL 3D
%%
\myslide{Photo d'espace-temps}%
{%
  \mySfigure{full}{dump}{anneal-3d}{\arule{Anneal} \acube{512}}}
{%
  ANNEAL 3D

  Une autre solution consiste \`a r\'eduire la quantit\'e de donn\'ees
  visualis\'es par des transformations implicitement associ\'ees \`a une
  technique de visualisation, ou explicitement appliqu\'e sur un
  espace-temps avant visualisation.

  Sur cette image la r\'eduction est implicite.
  
  La modulation des param\`etres d'une technique de visualisation permet
  de reconstruire la dimension temporelle.

  J'ais r\'ealis\'e des films de photos d'espace-temps en variant les
  param\`etre de projections.}

%%
%% QMEAN PIXMAP
%%
\myslide{Visualisation par pixmap}
{%
  \mySfigure{full}{dump}{tsum-grey}%
  {\arule{Qmean} \asquare{256} \agen{128} et \agen{256}}}
{%
  QMEAN PIXMAP

  Les g\'en\'erations 128 (\`a gauche) et 256 (\`a droite) d'une trajectoire
  de \arule{Qmean} sur un espace \asquare{256}.

  L'usage de pixmap pour la visualisation des configurations
  cellulaires de r\`egle non binaires pose le probl\`eme du choix de la
  colormap.

  Les colormaps sont ici des greymaps.

  Les deux pixmaps sup\'erieurs utilisent une greymap dont l'\'echelle
  des  valeurs est progressive.

  Les deux pixmaps inf\'erieurs utilisent une greymap dont l'\'echelle des
  valeurs est al\'eatoire, masquant la r\'epartition g\'en\'erale des
  intensit\'es, mais r\'ev\'elant les d\'etails de l'alternance de valeurs.}

%%%%%%%%%%%%%%%%%%%%%%%%%%%%%%%%%%%%%%%%%%%%%%%%%%%%%%%%%%%%%%%%

%%
%% SPL
%%
\begin{slide}
  \begin{Huge}
    \begin{center}
      SPL
    \end{center}
  \end{Huge}
\end{slide}

%%
%% CA DEFINITION
%%
\begin{slide}
  \begin{center}
    Les automates cellulaires
  \end{center}
  sont des syst\`emes dynamiques virtuels,\\
  discrets en espace, en temps et en \'etats,\\
  homog\`enes, synchrones et d\'eterministes.
\end{slide}

%%
%% THE RULES WE WILL USE
%%
\begin{slide}
  \mytitle{Les trois r\`egles utilis\'es\,:}
  \begin{tabbing}
    XXXXXXXXXXXX\=\kill
    \>\arule{Life}\\
    \>\arule{Anneal}\\
    \>\arule{Qmean}
  \end{tabbing}
\end{slide}

%%
%% SCHEMA FOR TRANSPARENCY MAKING
%%
\myslide{Sch\'ema de r\'ealisation de r\'eductions par transparence}
{\mySfigure{full}{fig}{transcub}{foo}}
{}

%%
%% SCHEMAS FOR SPACE-TIME SLICING
%%
\myslide{Projection }
{\mySfigure{full}{fig}{timepyrcub}{foo}}
{Tranches d'espace-temps et Pyramides relativistes}


%%
%% ANNEAL SPECIAL SECTION
%%

\begin{slide}
  \mystitle{Leur utilisation pour l'exploration}
  \mySfigure{full}{plot}{anneal-01a}%
  {\'Evolution de la taille des fronti\`eres de~\arule{Anneal}}%
  \mySfigure{full}{plot}{anneal-02}%
  {\'Evolution de la population de~\arule{Anneal} lors
    d'un changement du nombre d'\^iles}%
\end{slide}

%%\begin{mynote}
%%  La taille des fronti\`eres est une approximation par la diff\'erence bit
%%  \`a bit entre l'espace de g\'en\'eration~$n$ et~$n-1$.
%%
%%  Cette figure permet de comprendre l'existence des points
%%  d'inflexions de la courbe d'\'evolution de anneal.
%%\end{mynote}

%%
%%
%%

\begin{slide}
  \mySfigure{full}{plot}{anneal-03}%
  {Transition vers la derni\`ere \^ile}%
  \mySfigure{full}{plot}{anneal-04}%
  {Populations normalis\'ees pour 4~espaces initiaux de tailles crois\-santes}%
\end{slide}

%%\begin{mynote}
%%  Une autre zone de transition sur les m\^emes donn\'ees que celle de la
%%  figure pr\'ecedente.
%%  
%%  R\'eductions int\'egrales normalis\'ees des populations des trajectoires
%%  associ\'ees \`a 4~espaces initiaux de tailles respectives:~$2^{128}$,
%%  $2^{256}$, $2^{512}$ et $2^{1024}$. %
%%  Le plot de gauche pr\'esente le logarithme du plot de droite.
%%\end{mynote}

%%
%% LIFE XOR POPULATION PLOT
%%
\myslide{\'Evolutions de populations}
{\mySfigure{full}{plot}{lifex}{\arule{life} \asquare{128} \agencnt{1280}}}
{}

%%
%% QMEAN RAY-TRACED MOVIE
%%
\begin{slide}
  \mySfigure{full}{dump}{tsum-ray-all}%
  {\arule{Qmean} \acube{128}}
\end{slide}

\begin{mynote}
  Les variations de param\`etres sur l'axe horizontal porte sur la prise
  de vues.
  
  Les variations de param\`etres sur l'axe vertical porte sur la valeur
  de seuillage utilis\'ee.
  
  La valeur de seuil varie de 0.3~\`a~0.8.
\end{mynote}

%%
%% LIFE AND ANNEAL XOR SEQUENCE
%%
\begin{slide}
  \mytitle{Visualisations\,:}
  \mySfigure{full}{dump}{life-time-xor}%
  {\arule{Life}}
  \mySfigure{full}{dump}{anneal-time-xor}%
  {\arule{Anneal}}
\end{slide}

%%
%% QMEAN PIXMAP
%%
\myslide{Visualisation par pixmap}
{%
  \mySfigure{full}{dump}{tsum-grey}%
  {\arule{Qmean} \asquare{256} \agen{128} et \agen{256}}}
{%
  QMEAN PIXMAP

  Extraction param\'etr\'e par une colormap.

  \'echelle  des  valeurs est progressive.

  \'echelle des  valeurs est al\'eatoire\\
  masque la r\'epartition g\'en\'erale des intensit\'es\\
  r\'ev\'ele les d\'etails de l'alternance de valeurs.}

%%
%% QMEAN BITMAPS
%%
\myslide{Bitmaps de 8~plans de 8~g\'en\'erations}
{%
  \mySfigure{full}{dump}{tsum}%
  {\arule{Qmean} \asquare{128} \agencnt{1024}}}%
{%
  QMEAN BITMAPS

  Double extraction:\\
  tranche de temps et trance d'\'etats cellulaire

  Temps de haut en bas

  g\'en\'erations $2^{3}$ \`a $2^{10}$\\
  premi\`ere g\'en\'eration = 8.

  Les 256~\'etats possibles de chaque cellule sont visualis\'es par
  8~bitmaps au lieu d'un pixmap.

  Les bitmaps sont pr\'esent\'es de gauche \`a droite pour les bits~0 \`a~7.}

%%
%% ANNEAL AND LIFE TRANSPARENCY
%%
\begin{slide}
  \mystitle{Transparence par plan de bits}
  \mySfigure{full}{dump}{anneal-256x3-trans-bits}%
  {\arule{Anneal} \acube{256}}
  \mySfigure{full}{dump}{life-256x3-trans-bits}%
  {\arule{Life} \acube{256}}
\end{slide}

%%\begin{mynote}
%%  ANNEAL AND LIFE TRANSPARENCY
%%  
%%  R\'eduction par Accumulation.
%%\end{mynote}

% Local Variables: 
% mode: latex
% TeX-master: t
% End: 
